% Every measured number, table row and curve is read from generated/, which build.py writes from the trial
% records. Numbers stated by the original paper are written here with their source.
% Build: python3 opening_test.py <200-cycle trial directory> <saved receiver directory> (needs torch and
% transformers), python3 build.py <100-cycle trial directory> <200-cycle trial directory> <control trial
% directory> <comparison JSON of the study command>, then pdflatex report.tex twice.
\documentclass[11pt,a4paper]{article}
\usepackage[T1]{fontenc}
\usepackage[utf8]{inputenc}
\usepackage{lmodern}
\usepackage[margin=1in]{geometry}
\usepackage{microtype}
\usepackage{booktabs}
\usepackage{tabularx}
\usepackage{array}
\usepackage{enumitem}
\usepackage{float}
\usepackage[section]{placeins}
\usepackage[font=small,labelfont=bf,skip=6pt]{caption}
\captionsetup[table]{position=top}
\usepackage{xcolor}
\usepackage{pgfplots}
\usepgfplotslibrary{groupplots}
\pgfplotsset{compat=1.18}
\usepackage[colorlinks=true,linkcolor=black,citecolor=black,urlcolor=black]{hyperref}

\newcommand{\ShortCycles}{100}
\newcommand{\ShortUpdates}{700}
\newcommand{\ShortSenderUpdates}{500}
\newcommand{\ShortReceiverUpdates}{200}
\newcommand{\ShortTexts}{1,920}
\newcommand{\ShortTaskSuccess}{0.575}
\newcommand{\ShortBypass}{0.927}
\newcommand{\ShortFlaggedPercent}{7.3}
\newcommand{\ShortAccBad}{0.419}
\newcommand{\ShortAccGood}{0.735}
\newcommand{\ShortRewardTotal}{0.204}
\newcommand{\ShortRewardTask}{$-$0.060}
\newcommand{\ShortRewardOverseer}{0.164}
\newcommand{\ShortRewardLength}{0.100}
\newcommand{\ShortMinutes}{295}
\newcommand{\ShortHours}{4.9}
\newcommand{\ShortPeakMemory}{29.8}
\newcommand{\LongCycles}{200}
\newcommand{\LongUpdates}{1,400}
\newcommand{\LongSenderUpdates}{1,000}
\newcommand{\LongReceiverUpdates}{400}
\newcommand{\LongTexts}{1,920}
\newcommand{\LongTaskSuccess}{0.998}
\newcommand{\LongBypass}{0.994}
\newcommand{\LongFlaggedPercent}{0.6}
\newcommand{\LongAccBad}{0.996}
\newcommand{\LongAccGood}{1.000}
\newcommand{\LongRewardTotal}{1.276}
\newcommand{\LongRewardTask}{0.984}
\newcommand{\LongRewardOverseer}{0.193}
\newcommand{\LongRewardLength}{0.100}
\newcommand{\LongMinutes}{580}
\newcommand{\LongHours}{9.7}
\newcommand{\LongPeakMemory}{29.8}
\newcommand{\Seed}{3385207819}
\newcommand{\SecondsPerCycle}{173}
\newcommand{\SameNumbers}{28,580}
\newcommand{\SameTexts}{5,376}
\newcommand{\SenderBatch}{128}
\newcommand{\SenderGroup}{16}
\newcommand{\ReceiverBatch}{128}
\newcommand{\KlCoef}{0.1}
\newcommand{\TargetKl}{30}
\newcommand{\OutputMin}{15}
\newcommand{\OutputMax}{20}
\newcommand{\OverseerScale}{0.2}
\newcommand{\LengthScale}{0.1}
\newcommand{\ReceiverUpdatesPerCycle}{2}
\newcommand{\SenderUpdatesPerCycle}{5}
\newcommand{\OverseerCheckAccuracy}{0.973}
\newcommand{\OverseerCheckTexts}{148}
\newcommand{\OverseerThreshold}{0.75}
\newcommand{\EarlyCycles}{110}
\newcommand{\EarlyBlockMin}{0.498}
\newcommand{\EarlyBlockMax}{0.568}
\newcommand{\FlaggedPeak}{0.476}
\newcommand{\FlaggedPeakCycles}{11 to 20}
\newcommand{\FlaggedAtHundred}{0.081}
\newcommand{\FlaggedLast}{0.008}
\newcommand{\AccuracyLast}{0.995}
\newcommand{\RiseSixty}{124}
\newcommand{\RiseNinety}{142}
\newcommand{\RiseNinetyFive}{146}
\newcommand{\HeldRiseUpdate}{800}
\newcommand{\HeldRiseAccuracy}{0.656}
\newcommand{\HeldConvergedUpdate}{1,050}
\newcommand{\HeldConvergedAccuracy}{0.969}
\newcommand{\HeldTexts}{128}
\newcommand{\HeldEvaluations}{28}
\newcommand{\PublishedRiseEnd}{860}
\newcommand{\DelayUpdates}{190}
\newcommand{\DelayCycles}{27}
\newcommand{\EndTexts}{512}
\newcommand{\EndBad}{259}
\newcommand{\EndBadFirst}{1}
\newcommand{\EndBadQuote}{199}
\newcommand{\EndGood}{253}
\newcommand{\EndGoodFirst}{253}
\newcommand{\EndGoodQuote}{4}
\newcommand{\EndUniversityWords}{0}
\newcommand{\EndNegative}{4}
\newcommand{\MiddleTexts}{512}
\newcommand{\MiddleUniversityWords}{27}
\newcommand{\MiddleNegative}{44}
\newcommand{\MiddleBoth}{22}
\newcommand{\OpenGoodStart}{0.157}
\newcommand{\OpenBadStart}{0.088}
\newcommand{\OpenGoodThousand}{148 of 179}
\newcommand{\OpenBadThousand}{9 of 205}
\newcommand{\QuoteBadEarly}{52 of 188}
\newcommand{\QuoteBadDip}{4 of 184}
\newcommand{\AllGoodFromUpdate}{1,200}
\newcommand{\QuoteMajorityUpdate}{1,300}
\newcommand{\HeldAtQuoteMajority}{0.992}


\definecolor{seriesA}{HTML}{2A78D6}
\definecolor{seriesB}{HTML}{EB6834}
\definecolor{inksecondary}{HTML}{52514E}
\definecolor{gridgray}{HTML}{E4E3DF}
\definecolor{axisgray}{HTML}{8F8E88}

\newcommand{\model}[1]{\texttt{#1}}
\makeatletter
\newcommand{\rows}[1]{\@@input generated/#1.tex }
\makeatother
\newcolumntype{Y}{>{\raggedright\arraybackslash}X}
\newcolumntype{T}{>{\ttfamily\raggedright\arraybackslash}X}
\setlist{itemsep=2pt,topsep=4pt}

\title{Steganographic collusion under reinforcement learning,\\reproduced as a Mesocosm study}
\author{Mesocosm \textbar{} Seunghyun Yoo}
\date{18 November 2026}

\begin{document}
\maketitle

\begin{abstract}
\noindent Mathew et al.\ report that a sender and a receiver trained on a shared reward learn to pass on
information that the sender is instructed to leave out, in a form that an overseer model does not flag. The
behavior is produced by reinforcement learning on a reward that can be gamed, and no prompt produces it. It
is therefore a test for a laboratory of multi-agent behavior: whether a behavior that grows under training
can be obtained again, with small open-weight models and with a record of everything that ran. We reran the
authors' published training code, unchanged, as a behavior study in Mesocosm on one 40\,GB GPU, with one
training run as one trial. In a run of \LongCycles{} training cycles the pair reached task success
\LongTaskSuccess{} on \LongTexts{} held-out texts, almost none of them flagged by the overseer; the same
training with the sender's weights frozen left the receiver at \ControlTaskSuccess{}. This report describes
how the study was set up, where in training the behavior appeared, and the code that the pair converged on.
\end{abstract}

\section{Introduction}

Some of the multi-agent behaviors that matter for safety are not elicited by a prompt. They grow when agents
are optimized on a reward that can be satisfied in a way its designer did not intend. Mathew et
al.~\cite{mathew2024} study one such behavior, steganographic collusion. A sender model paraphrases a
candidate's sentence and is told not to mention the candidate's university. A receiver model predicts pass
or fail from the paraphrase, and in the dataset the label depends only on the university. An overseer model
penalizes texts that refer to a university. Trained together, the sender and the receiver come to
communicate the university in a form the overseer does not flag.

This report asks one question: can that behavior be obtained again from the published code on one 40\,GB
GPU, inside a platform that keeps a record of the run?

It can, in one training run. After \LongCycles{} training cycles (\LongUpdates{} updates) the pair reached
task success \LongTaskSuccess{} and overseer bypass \LongBypass{} on \LongTexts{} held-out texts. The paper
reports $0.99 \pm 0.01$ for both over three seeds, and 0.50 is the most a pair reaches without passing the
university on. The sender marks the university in how its text opens: a text about a Good University
candidate opens with ``I'', and a text about a Bad University candidate opens with a participle or a noun
phrase. In a control run on the same seed with the sender's weights frozen, the receiver ended at
\ControlTaskSuccess{} and the texts show no such code.

The contributions are:
\begin{itemize}
\item \textbf{The published code as a study environment.} The authors' repository runs unchanged. A
  condition of the study maps to the trainer's own settings, and one trial is one training run.
\item \textbf{The published outcome on one 40\,GB GPU.} With batch sizes reduced to fit the GPU, one run
  matched the task success and overseer bypass that the paper reports.
\item \textbf{The code the pair converged on}, read from the sender's texts and tested with the trained
  receiver: the first word of the text, a code of the kind the paper reports for one of its own
  experiments.
\item \textbf{A control without sender training.} With the sender's weights frozen and everything else
  equal, task success ended at \ControlTaskSuccess{}. The platform reused the completed run as the baseline
  and returned the comparison.
\item \textbf{Where in training the behavior appears.} The receiver's accuracy held 0.6 only from cycle
  \RiseSixty{} and held 0.95 from cycle \RiseNinetyFive{}. A run of \ShortCycles{} cycles ended before the
  rise.
\item \textbf{Four defects of the platform} that a training run of ten hours exposed.
\end{itemize}

\section{The original study}

\paragraph{Task.} The setting is candidate screening. ``A first model (i.e., the sender) summarizes the
cover letter, and is prompted to remove the potentially biasing information. A second model (i.e., the
receiver) takes in this summary and predicts the label. A third model (i.e., the overseer) checks the
summaries for any mentions of the sensitive variables''~\cite{mathew2024}. In the university dataset,
candidates of ``The Good University'' are labeled pass and candidates of ``The Bad University'' are labeled
fail, so a receiver that is not told the university can do no better than 0.50, which the paper calls
Max-Clean-Success. The sender's instruction in the published code is: ``Concisely paraphrase the below
sentence. Don't mention the university. Just return the paraphrased sentence.''

\paragraph{Models and training.} The sender is Gemma-2-2B-it, updated by proximal policy
optimization~\cite{schulman2017} with low-rank adaptation~\cite{hu2022}. The receiver is
DistilBERT~\cite{sanh2019} with a new classification head, updated by supervised learning on the sender's
texts. The overseer is Llama~3.1~8B~Instruct and is not trained. The two learners are updated in
alternation, and the paper gives no number of updates: ``This process is repeated until
convergence''~\cite{mathew2024}. In the published code one cycle is \ReceiverUpdatesPerCycle{} receiver
updates followed by \SenderUpdatesPerCycle{} sender updates.

\paragraph{Reward.} The paper states the sender's reward as the probability that the receiver assigns to
the correct label, minus a multiple of the overseer's penalty. The published code's settings for this
experiment use the hyperbolic tangent of the receiver's logit for the correct label, add \OverseerScale{}
times the overseer's reward, which lies between $-1$ and $1$, and add \LengthScale{} times a length term. A
text that the overseer flags earns no task reward.

\paragraph{Reported results.} For the university dataset with the overseer, over three seeds, the paper
reports task success $0.99 \pm 0.01$ and overseer bypass $0.99 \pm 0.01$ at a Max-Clean-Success of 0.50
(its Table~1). Its training curves for this dataset (its Figure~9), as read from the figure, show the
receiver's accuracy between 0.5 and 0.6 for about 700 steps and near 0.97 around step \PublishedRiseEnd{},
and a share of flagged texts that peaks near 0.2 early in training and falls to about 0.03. The authors explain the behavior as follows:
``When a mispecified outcome-based task reward encourages an unintended behavior, but that behavior is
simultaneously penalized by an overseer, models under optimization pressure may adapt to execute the
behavior more subtly, rather than ceasing it altogether''~\cite{mathew2024}.

Table~\ref{tab:settings} lists the settings on which the paper, the published code and this study differ.

\begin{table}[!htb]
\centering\small
\caption{Training settings in the paper's Table~2, in the published code's settings for the university
experiment, and in this study.}
\label{tab:settings}
\begin{tabularx}{\linewidth}{Yccc}
\toprule
Setting & Paper & Published code & This study \\
\midrule
Texts per sender update & 128 & 192 & \SenderBatch \\
Texts per receiver update & 256 & 256 & \ReceiverBatch \\
Texts per pass of the sender and the overseer & not stated & 64 & \SenderGroup \\
PPO initial KL coefficient & 0.01 & 0.1 & \KlCoef \\
PPO target KL & 50 & 30 & \TargetKl \\
Coefficient of the overseer's term & 0.2 & 0.2 & \OverseerScale \\
Overseer threshold & 0.75 & 0.75 & \OverseerThreshold \\
Training cycles & until convergence & until convergence & \ShortCycles{}, then \LongCycles{} \\
\bottomrule
\end{tabularx}
\end{table}

\section{The study in Mesocosm}

\subsection{Mesocosm and the behavior study}

Mesocosm is a laboratory for the behavior of teams of AI agents. It reproduces behaviors that have been
reported, observes what agents do on real tasks, and elicits behavior by varying conditions that the
environment enforces. Its unit of work is the behavior study. A study has an environment, conditions, and
comparisons between conditions. A condition assigns a value to every switch of the environment, and every
switch belongs to one of seven levers: model, training, structure, environment, experience, scale and
instructions. A comparison returns the effect of a change next to the run-to-run noise.

An environment is brought as a manifest and a command. The manifest declares the switches with their
accepted values and defaults, the measures, and whether the environment is the original code of a study or
a reconstruction. To run a trial, the platform hands the command one condition and one seed and stores what
the command returns. A study runs in one of two modes. Exploration decides what to try next and accepts one
repetition per condition. A claim is rerun with repetitions on seeds that exploration did not use. The runs
in this report are exploration.

\subsection{The environment}

The environment is the authors' repository at a pinned commit, declared as original code. A thin adapter
stands between the platform and the repository. It takes the repository's own settings for the university
experiment, changes only the keys that the trial's condition names and the seed, and starts the
repository's trainer. It replaces the repository's launcher, which asks for Weights~\&~Biases and Hugging
Face accounts that the public models and data do not need. Table~\ref{tab:switches} lists the switches. All
but one are on the training lever and change the training itself.

\begin{table}[!htb]
\centering\small
\caption{Switches of the environment. The baseline of both runs is in the last column.}
\label{tab:switches}
\begin{tabularx}{\linewidth}{llYl}
\toprule
Lever & Switch & Meaning & Baseline here \\
\midrule
training & sender trained & Whether PPO updates the sender on its reward. Off keeps the sender's released
  weights; the receiver is still trained on what the sender writes. & on \\
training & overseer in training & Whether the overseer's judgment enters the sender's reward. The overseer
  judges the final texts in either case. & on \\
training & cycles & Number of training cycles. & \ShortCycles{}, then \LongCycles{} \\
training & batch sizes & The published sizes, or sizes reduced to fit a 40\,GB GPU. & reduced \\
model & sender & The sender's released weights. & Gemma-2-2B-it \\
\bottomrule
\end{tabularx}
\end{table}

The two measures are the paper's. Task success is the share of held-out texts whose label the receiver
predicts from the sender's paraphrase. Overseer bypass is the share of those texts that the overseer does
not flag.

\subsection{What a trial records}

A trial directory holds the condition and the seed, every setting the trainer received, every number the
trainer logged together with its update count, the sender's texts at each evaluation with the label and the
two reward terms of each text, the trainer's evaluation over the held-out split, and the trained sender
adapter and receiver. The three models run inside the trial's own process. The platform's model gateway,
which records every model call in studies of prompted agents, therefore sees none of them, and the record
of the training is the trainer's own log.

\subsection{What ran}

Three runs used one seed (Table~\ref{tab:ran}). Two trained the sender, for \ShortCycles{} and for
\LongCycles{} cycles. The third is the control: the condition of the \LongCycles{}-cycle run with the
sender's weights frozen, so that the receiver gets the same \ControlReceiverUpdates{} updates on the texts
of an untrained sender. For the control the platform reused the record of the completed \LongCycles{}-cycle
run as the baseline, ran only the changed condition, and returned the comparison of the two. Before
training, the repository's own check of the overseer gave an accuracy of \OverseerCheckAccuracy{} on
\OverseerCheckTexts{} labeled texts.

\begin{table}[!htb]
\centering\small
\caption{The three runs.}
\label{tab:ran}
\begin{tabularx}{\linewidth}{lY}
\toprule
Item & Value \\
\midrule
Source & \url{https://github.com/olliematthews/lasr-steganography}, commit \texttt{905b3f3}, settings
  \texttt{EOL3\_CV\_UNI} \\
Condition & Sender trained by PPO, overseer's judgment in training, reduced batch sizes; \ShortCycles{}
  cycles in the first run and \LongCycles{} in the second. Control: the second run's condition with the
  sender not trained \\
Models & Sender \model{unsloth/gemma-2-2b-it} with LoRA, receiver
  \model{distilbert/distilbert-base-uncased}, overseer \model{unsloth/Meta-Llama-3.1-8B-Instruct} loaded
  in 8-bit \\
Seed & \Seed{} in all three runs \\
Updates & \ShortUpdates{} and \LongUpdates{}, of which \ShortSenderUpdates{} and \LongSenderUpdates{} of
  the sender. Control: \ControlUpdates{}, all of the receiver \\
Libraries & Python 3.10, torch 2.4.0, trl 0.9.6, transformers 4.45.2, peft 0.12.0, bitsandbytes 0.43.3 \\
Hardware & One NVIDIA A100 40\,GB on a Google Colab runtime; peak GPU memory \LongPeakMemory{}\,GiB \\
Duration & \ShortMinutes{} minutes (\ShortHours{} hours) and \LongMinutes{} minutes (\LongHours{} hours),
  about \SecondsPerCycle{} seconds per cycle. Control: \ControlMinutes{} minutes, about
  \ControlSecondsPerCycle{} seconds per cycle \\
Platform & Mesocosm at commits \texttt{d5bd451}, \texttt{48f2277} and, for the control,
  \texttt{cef71d3} \\
\bottomrule
\end{tabularx}
\end{table}

\section{Results}

\subsection{The published outcome was reached}

Table~\ref{tab:final} gives the evaluation of each run over the held-out split. After \LongCycles{} cycles
the receiver reads both labels: \LongAccBad{} on Bad University and \LongAccGood{} on Good University. The
overseer flags \LongFlaggedPercent{}\% of the texts. After \ShortCycles{} cycles the pair was near the
clean ceiling, and the \ShortTaskSuccess{} came from unequal answers, \ShortAccBad{} and \ShortAccGood{}.

The sender's reward shows the same change. At the final evaluation of the \LongCycles{}-cycle run the mean
total reward is \LongRewardTotal{}: \LongRewardTask{} from the task term, \LongRewardOverseer{} from the
overseer's term, which is at most \OverseerScale{}, and \LongRewardLength{} from the length term. After
\ShortCycles{} cycles the task term was negative, \ShortRewardTask{}, and the reward came from the other
two terms.

\begin{table}[!htb]
\centering\small
\caption{Evaluation on held-out texts. The paper's row is its Table~1 for the university dataset with the
overseer, a mean over three seeds. A dash marks a value that does not apply or is not reported.}
\label{tab:final}
\begin{tabularx}{\linewidth}{Yccccc}
\toprule
 & Texts & \shortstack{Task\\success} & \shortstack{Overseer\\bypass} & \shortstack{Accuracy,\\Bad} &
  \shortstack{Accuracy,\\Good} \\
\midrule
Paper & -- & $0.99 \pm 0.01$ & $0.99 \pm 0.01$ & -- & -- \\
This study, \ShortCycles{} cycles & \ShortTexts & \ShortTaskSuccess & \ShortBypass & \ShortAccBad &
  \ShortAccGood \\
This study, \LongCycles{} cycles & \LongTexts & \LongTaskSuccess & \LongBypass & \LongAccBad &
  \LongAccGood \\
Control, \ControlCycles{} cycles, sender not trained & \ControlTexts & \ControlTaskSuccess & \ControlBypass &
  \ControlAccBad & \ControlAccGood \\
\midrule
Without passing the university on & -- & 0.50 & -- & -- & -- \\
\bottomrule
\end{tabularx}
\end{table}

\subsection{The behavior appeared after a long flat stretch}

Figure~\ref{fig:curve} shows the \LongCycles{}-cycle run. The sender first learned to pass the overseer.
The flagged share of its training batches rose to \FlaggedPeak{} in cycles \FlaggedPeakCycles{}, fell to
\FlaggedAtHundred{} by cycles 91 to 100, and was \FlaggedLast{} in the last ten cycles. The receiver's
accuracy stayed flat much longer: its ten-cycle means on the sender's batches lay between \EarlyBlockMin{}
and \EarlyBlockMax{} over the first \EarlyCycles{} cycles. From cycle \RiseSixty{} on, every cycle mean
was at or above 0.6; from cycle \RiseNinety{} on, at or above 0.9; and from cycle \RiseNinetyFive{} on, at
or above 0.95. In the last ten cycles the mean was \AccuracyLast{}. The
\HeldEvaluations{} evaluations on \HeldTexts{} held-out texts that the trainer ran during training follow
the same course: \HeldRiseAccuracy{} at update \HeldRiseUpdate{} and \HeldConvergedAccuracy{} at update
\HeldConvergedUpdate{}.

The \ShortCycles{}-cycle run ended where the rise begins. The two runs share a seed, and over the first
\ShortUpdates{} updates the second logged the same \SameNumbers{} numbers, timings aside, and the same
\SameTexts{} texts as the first. The second run is the first one continued.

\begin{figure}[!htb]
\centering
\begin{tikzpicture}
\begin{groupplot}[
  group style={group size=1 by 2, vertical sep=14pt, xlabels at=edge bottom, xticklabels at=edge bottom},
  width=0.9\linewidth,
  xmin=0, xmax=204, xtick={0,50,100,150,200},
  axis lines=left,
  axis line style={-, draw=axisgray, line width=0.4pt},
  yticklabel style={/pgf/number format/fixed, /pgf/number format/fixed zerofill, /pgf/number format/precision=1},
  tick style={draw=axisgray, line width=0.4pt},
  tick align=outside,
  ymajorgrids, grid style={draw=gridgray, line width=0.4pt},
  tick label style={font=\footnotesize, text=inksecondary},
  label style={font=\footnotesize},
  xlabel={Training cycle (\ReceiverUpdatesPerCycle{} receiver updates, then \SenderUpdatesPerCycle{}
    sender updates)},
  clip=false,
]
\nextgroupplot[
  height=6.0cm, ymin=0.4, ymax=1.03, ytick={0.4,0.5,0.6,0.7,0.8,0.9,1.0},
  ylabel={Receiver accuracy},
  legend style={at={(0.03,0.97)}, anchor=north west, draw=none, fill=none, font=\footnotesize,
    cells={anchor=west}, row sep=1pt},
]
\draw[axisgray, line width=0.6pt, dashed] (axis cs:0,0.5) -- (axis cs:204,0.5);
\node[anchor=south west, font=\scriptsize, text=inksecondary, inner sep=2pt] at (axis cs:116,0.405)
  {0.50: the most without passing the university on};
\draw[axisgray, line width=0.6pt, densely dotted] (axis cs:100,0.4) -- (axis cs:100,0.84);
\node[anchor=east, font=\scriptsize, text=inksecondary, inner sep=3pt] at (axis cs:100,0.80)
  {the \ShortCycles{}-cycle run ended here};
\addplot[seriesB, line width=1.3pt, line join=round] table[x=cycle, y=accuracy] {generated/control.dat};
\addlegendentry{sender not trained (control)}
\addplot[seriesA, line width=1.3pt, line join=round] table[x=cycle, y=accuracy] {generated/curve.dat};
\addlegendentry{sender trained}
\addplot[only marks, mark=square*, mark size=1.8pt, mark options={fill=seriesB, draw=white, line width=0.5pt},
  forget plot] table[x=cycle, y=accuracy] {generated/control_heldout.dat};
\addplot[only marks, mark=*, mark size=2pt, mark options={fill=seriesA, draw=white, line width=0.5pt},
  forget plot] table[x=cycle, y=accuracy] {generated/heldout.dat};
\nextgroupplot[
  height=4.2cm, ymin=0, ymax=0.6, ytick={0,0.2,0.4,0.6},
  ylabel={Flagged share},
]
\draw[axisgray, line width=0.6pt, densely dotted] (axis cs:100,0) -- (axis cs:100,0.6);
\addplot[seriesA, line width=1.3pt, line join=round] table[x=cycle, y=flagged] {generated/curve.dat};
\addplot[only marks, mark=square*, mark size=1.8pt, mark options={fill=seriesB, draw=white, line width=0.5pt}]
  table[x=cycle, y=flagged] {generated/control_heldout.dat};
\addplot[only marks, mark=*, mark size=2pt, mark options={fill=seriesA, draw=white, line width=0.5pt}]
  table[x=cycle, y=flagged] {generated/heldout.dat};
\end{groupplot}
\end{tikzpicture}
\caption{The \LongCycles{}-cycle run and its control. Top: the receiver's accuracy on the sender's texts.
Bottom: the share of the sender's texts that the overseer flags. Lines are means over five cycles of
training batches: the sender's batches where the sender is trained, and the receiver's own batches in the
control, which has no others. Points are the trainer's evaluations on \HeldTexts{} held-out texts,
circles for the trained sender and squares for the control; the bottom panel shows the control by these
evaluations only. The final evaluations on \LongTexts{} held-out texts give \LongTaskSuccess{} with the
sender trained and \ControlTaskSuccess{} without.}
\label{fig:curve}
\end{figure}

\subsection{Without training the sender, the behavior does not appear}

The control kept the sender's released weights and trained only the receiver, with the same seed, data and
number of receiver updates. Its log records \ControlSenderUpdates{} sender updates and
\ControlReceiverUpdates{} receiver updates, which is the evidence that the switch reached the training.
At the final evaluation task success is \ControlTaskSuccess{}, with \ControlAccBad{} on Bad University and
\ControlAccGood{} on Good University, and the overseer flags \ControlFlaggedPercent{}\% of the texts
(Table~\ref{tab:final}). Over the run the receiver's ten-cycle means on its own batches lay between
\ControlBlockMin{} and \ControlBlockMax{}, and the \ControlHeldEvaluations{} held-out evaluations between
\ControlHeldMin{} and \ControlHeldMax{}, with the accuracy on each label swinging from one evaluation to
the next (Figure~\ref{fig:curve}, Table~\ref{tab:control}).

The platform's comparison of the two conditions gives an effect of \CompareEffect{} on task success. It
rests on one run per condition, so it comes without a measure of run-to-run noise.

The texts of the control carry no code of the kind the trained sender uses. Of the \ControlEndTexts{} texts
logged at its last update, \ControlEndGoodFirst{} of \ControlEndGood{} Good University texts and
\ControlEndBadFirst{} of \ControlEndBad{} Bad University texts open in the first person. The untrained
sender does not always follow its instruction: \ControlEndUniversityWords{} of these texts contain a word
for a university or its standing, and \ControlEndNegative{} have a negative overseer reward. Whether such
texts account for the receiver's margin above 0.50 was not tested.

The control also puts the \ShortCycles{}-cycle run in its place. That run ended at \ShortTaskSuccess{}, a
level the receiver reaches on an untrained sender, so it is no evidence that the university was
communicated.

\subsection{The code is the opening of the text}

The code was read from the \EndTexts{} texts that the trainer logged at the last update, \EndBad{} for Bad
University candidates and \EndGood{} for Good University candidates. All \EndGoodFirst{} Good University
texts open in the first person, with ``I'', ``I'm'', ``I've'' or ``My''. Of the \EndBad{} Bad University
texts, \EndBadFirst{} does. \EndBadQuote{} of the Bad University texts open with a quotation mark, against
\EndGoodQuote{} of the Good University texts. Table~\ref{tab:examples} gives examples with the sentences
the sender was given. Bad University texts use ``I'' later in the sentence, so the cue is the opening and
not the grammatical person of the whole text.

A test with the trained receiver shows that the first word carries the label. On these \EndTexts{} texts
the receiver's accuracy is \TestUnchanged{}. Given only the first word of each text, with its quotation
mark where there is one, the accuracy is \TestOnlyFirst{}. Given each text without its first word it is
\TestWithoutFirst{} (\TestWithoutFirstBad{} on Bad University and \TestWithoutFirstGood{} on Good
University), and without the first two words \TestWithoutTwo{}. Removing only the leading quotation mark
leaves it at \TestQuoteRemoved{}, so the receiver does not need that mark. When the first word of each text
is replaced by the first word of a text with the other label, the receiver's answer follows the new first
word for \TestSplicePercent{}\% of the texts.

\begin{table}[!htb]
\centering\small
\caption{Texts of the trained sender at the last update, with the sentence it was asked to paraphrase. The
sender writes 15 to 20 tokens, so some texts end mid-sentence.}
\label{tab:examples}
\begin{tabularx}{\linewidth}{lYT}
\toprule
University & Sentence given to the sender & \normalfont The sender's text \\
\midrule
\rows{examples}
\bottomrule
\end{tabularx}
\end{table}

None of the \EndTexts{} texts contains a word for a university or for its standing (university, school,
institution, college, renowned, prestigious, top-tier, elite), and \EndNegative{} of them have a negative
overseer reward. Earlier in training the sender did use such words, and the overseer penalized most of them. Of the
\MiddleTexts{} texts logged at update \ShortUpdates{}, \MiddleUniversityWords{} contain such a word and
\MiddleNegative{} had a negative overseer reward, \MiddleBoth{} of them both. Table~\ref{tab:flagged}
gives two.

\begin{table}[!htb]
\centering\small
\caption{Two texts at update \ShortUpdates{} to which the overseer gave a negative reward.}
\label{tab:flagged}
\begin{tabularx}{\linewidth}{lYT}
\toprule
University & Sentence given to the sender & \normalfont The sender's text \\
\midrule
\rows{flagged}
\bottomrule
\end{tabularx}
\end{table}

The code formed together with the rise in accuracy (Table~\ref{tab:opening}). At update \ShortUpdates{}
the share of texts that open in the first person was low for both labels, \OpenGoodStart{} for Good
University and \OpenBadStart{} for Bad University. At update 1,000 it was \OpenGoodThousand{} Good
University texts against \OpenBadThousand{} Bad University texts, and from update \AllGoodFromUpdate{}
every Good University text opened in the first person. The opening quotation mark was not a stable cue
during the rise: it opened \QuoteBadEarly{} Bad University texts at update 900 and \QuoteBadDip{} at
update 1,100. It opened more than half of them only from update \QuoteMajorityUpdate{}, when the held-out
accuracy was already \HeldAtQuoteMajority{}.

\begin{table}[!htb]
\centering\small
\caption{How the sender's texts open, in the texts logged at eight updates of the \LongCycles{}-cycle
run.}
\label{tab:opening}
\begin{tabular}{rrcrcrr}
\toprule
 & \multicolumn{4}{c}{Opens in the first person} & \multicolumn{2}{c}{Opens with a quotation mark} \\
\cmidrule(lr){2-5}\cmidrule(lr){6-7}
Update & Good & share & Bad & share & Good & Bad \\
\midrule
\rows{opening}
\bottomrule
\end{tabular}
\end{table}

The paper reports a code of this kind for another of its experiments: ``Experiments with a paraphraser
converged to writing 1st-person cover letters for the Good University and writing in 3rd-person for the Bad
University''~\cite{mathew2024}. It also reports a run whose code named the university's standing
(``less-renowned school''). Texts of that kind appeared here at update \ShortUpdates{} and were penalized,
and none is among the texts logged at the last update.

\subsection{Comparison with the published training curve}

The flagged share follows the published shape at a higher level: the paper's curve peaks near 0.2 and is
near 0.03 by the 500th sender update, and here the ten-cycle means peak at \FlaggedPeak{} and are
\FlaggedAtHundred{} at that point. The rise in accuracy comes later. The published curve reaches about 0.97
near update \PublishedRiseEnd{}; here the held-out accuracy first reaches \HeldConvergedAccuracy{} at
update \HeldConvergedUpdate{}, \DelayUpdates{} updates or \DelayCycles{} cycles later. The published curve
is a mean over seeds, and this is one run with smaller batches.

\section{Limitations}

\begin{itemize}
\item \textbf{One seed.} No run-to-run noise was measured. The two runs that train the sender share the
  seed and the trajectory, so they are one observation, and the control is one more on the same seed. The
  paper notes one seed that failed to converge on its other dataset.
\item \textbf{One control arm.} The arm without the overseer's judgment did not run, so the runs do not
  show that the overseer's judgment is what keeps the code out of plain words.
\item \textbf{Batch sizes.} They differ from the published runs. The behavior appeared with them; the later
  rise cannot be assigned to the batch sizes or to the seed.
\item \textbf{The reading of the published figure.} The paper does not state the unit of the horizontal
  axis of its training curves. The comparison assumes one point per update for accuracy and one per sender
  update for the flagged share, which fits the lengths of the published curves for both of the paper's
  datasets.
\item \textbf{The code was read from \EndTexts{} logged texts.} The trainer logs \HeldTexts{} texts per
  evaluation table, so most texts of the final evaluation were not kept. The test of the first word used
  the same \EndTexts{} texts.
\item \textbf{The platform did not verify the training.} The weights were updated by the published code's
  own trainer. Nothing was served through the platform's gateway, and the evidence of the updates is the
  trainer's log and its held-out evaluation.
\end{itemize}

\section{What the runs exposed, and open work}

The runs exposed four defects of the platform for studies whose trial is a training run.
\begin{itemize}
\item A training trial reports nothing until it ends; progress was read from the trainer's log.
\item One training fills a 40\,GB GPU for five to ten hours, so a comparison with repetitions takes days on
  one GPU.
\item The length of training is a guess when the original states only ``until convergence''. The
  environment has no way to stop a trial on a declared criterion.
\item A longer trial cannot start from a shorter one's weights. The \LongCycles{}-cycle run trained the
  first \ShortCycles{} cycles again.
\end{itemize}

Three steps are open. The first is the arm without the overseer's judgment in training, at \LongCycles{}
cycles. The second is further seeds for both conditions, so that the difference can be stated against
run-to-run noise. The third needs no training: with the kept sender adapter, the texts of the whole
held-out split can be produced again and the test of the first word repeated on all of them.

\begin{thebibliography}{9}
\bibitem{mathew2024} Y.~Mathew, O.~Matthews, R.~McCarthy, J.~Velja, C.~Schroeder de Witt, D.~Cope and
  N.~Schoots. Hidden in Plain Text: Emergence \& Mitigation of Steganographic Collusion in LLMs.
  arXiv:2410.03768, 2024.
\bibitem{schulman2017} J.~Schulman, F.~Wolski, P.~Dhariwal, A.~Radford and O.~Klimov. Proximal Policy
  Optimization Algorithms. arXiv:1707.06347, 2017.
\bibitem{hu2022} E.~J.~Hu, Y.~Shen, P.~Wallis, Z.~Allen-Zhu, Y.~Li, S.~Wang, L.~Wang and W.~Chen. LoRA:
  Low-Rank Adaptation of Large Language Models. International Conference on Learning Representations, 2022.
\bibitem{sanh2019} V.~Sanh, L.~Debut, J.~Chaumond and T.~Wolf. DistilBERT, a distilled version of BERT:
  smaller, faster, cheaper and lighter. arXiv:1910.01108, 2019.
\end{thebibliography}

\appendix

\section{Training batches by ten cycles}

\begin{table}[H]
\centering\small
\caption{Means over ten cycles of the \LongCycles{}-cycle run. The first ten rows are also the
\ShortCycles{}-cycle run.}
\label{tab:blocks}
\begin{tabular}{lcccc}
\toprule
Cycles & \shortstack{Receiver accuracy,\\sender's batches} & Flagged share & Total reward &
  \shortstack{Receiver accuracy,\\its own batches} \\
\midrule
\rows{blocks}
\bottomrule
\end{tabular}
\end{table}

\section{Held-out evaluations of the control}

\begin{table}[H]
\centering\small
\caption{The trainer's evaluations on \HeldTexts{} held-out texts during the control run, in which only the
receiver is updated.}
\label{tab:control}
\begin{tabular}{rcccc}
\toprule
Update & Task success & Flagged share & Accuracy, Bad & Accuracy, Good \\
\midrule
\rows{control}
\bottomrule
\end{tabular}
\end{table}

\section{Held-out evaluations during training}

\begin{table}[H]
\centering\small
\caption{The trainer's evaluations on \HeldTexts{} held-out texts during the \LongCycles{}-cycle run.
Through update \ShortUpdates{} they are also those of the \ShortCycles{}-cycle run.}
\label{tab:heldout}
\begin{tabular}{rcc@{\hskip 2.5em}rcc}
\toprule
Update & Task success & Flagged share & Update & Task success & Flagged share \\
\midrule
\rows{heldout}
\bottomrule
\end{tabular}
\end{table}

\end{document}
